\section{Virtual and Real Metaphyics}

Metaphysics is an important element of all Traditions. By that we mean an exact science, as precise as mathematics, and as such, it is neither a speculative system nor a matter for philosophical debate. The fundamental tenet of metaphysics is that there are principles that determine or explain the hidden structure of the world. These principles can be understood either in a virtual or a real sense.

\begin{itemize}
\item \textbf{Virtual}. Guenon explains that metaphysics is beyond logic, but is not illogical. That is, there is no compelling logical argument to coerce a thinking person to accept the first principles it proposes. Nevertheless, given those first principles, other principles do follow logically from them. Someone may accept these principles as a matter of thought; this is virtual understanding. 
\item \textbf{Real}. Beyond the merely verbal understanding, Guenon refers to a ``metaphysical realization". That is, the principles are known in a direct way, or by an ``intuition", in analogy to the direct perception of physical objects. In this realization, there is a change in state of being to correspond to the principle. Intuition, in this sense, does not mean anything like a ``hunch" or any other irrational form of knowing. 
\end{itemize}
The first step is a virtual understanding. Based on the testimonies of the seers and saints of tradition regarding real understanding, a man has faith in the truth of the principles. Through initiation and spiritual practice, that understanding deepens. Since man is a microcosm, these practices involve dispassionate observation of one's own states of consciousness; this understanding then is also an understanding of the macrocosm.

Post-traditional systems of thought involve the rejection of these metaphysical principles. In order to understand those systems from within, it is neither unreasonable not unintelligent to determine which principles are rejected. To the contrary, the most intelligent thinkers know explicitly what they reject. Their followers, on the other hand, take a system at face value, not bothering to understand the premises on which it rests. Despite the many variations of such philosophical systems over the course of time, there are simply not many sensible alternatives.

For example, among the systems that reject virtual understanding, these are some options:

\begin{itemize}
\item \textbf{Empiricism}. This rejects out of hand the very idea of first principles. Knowledge can only be achieved by observing the physical world and forming scientific theories. Knowledge thus accumulates over time, but can never be complete. The civilizations of the past were defective because they lack the latest scientific knowledge. 
\item \textbf{Materialism}. The first principles of metaphysics may be replaced by others, such as materialism, determinism, and related ideas. Obviously, this denies intelligibility to the world, since everything happens only through material processes. Man, himself, is a material object, a being randomly thrown into the world. Hence, man as a person does not really exist. In actuality, this position can never be held consistently since the conscious awareness of intelligent choices and will are too strong. 
\item \textbf{Voluntarism}. In this case, the first principle is the Will, as in the systems developed by Schopenhauer and Nietzsche. Once again, if Will is primary, then Intelligence must be derivative. These are reactions to Kant since, if the pure reason can never grasp ultimate reality, then the practical reason, or Will, is the only way to interact with the world. Unfortunately, this also denies the ultimate intelligibility of the world, although it aligns more with personal experience than does materialism. 
\end{itemize}
Some of the consequences of denying the real understanding of metaphysics are these:

\begin{itemize}
\item \textbf{Fideism}. Fideism does not deny transcendence, but it denies that it can be known. Instead, it can only be a matter of ``faith". This equates faith with credulity, a lesser thing than the Traditional understanding of faith. 
\item \textbf{Naturalism}. These systems deny transcendence itself. Only the physical world exists and any claims to the contrary are regarded as the products of superstition and ignorance. 
\end{itemize}
This list may be incomplete and philosophical systems may mix some of these together. This is not to say that the conclusions of such thinkers are uninteresting, unintelligent, or untrue. Rather, the point is that it is futile to debate at that level since fundamental presuppositions are incompatible and thus preclude the possibility of a common understanding. To avoid succumbing to modern thought, it is not only useful, but also necessary, to bring these presuppositions to light.



\flrightit{Posted on 2012-09-25 by Cologero }

\begin{center}* * *\end{center}

\begin{footnotesize}\begin{sffamily}



\texttt{Cologero on 2012-09-26 at 07:34 said: }

As an example of inconsistent thinking, we can point to these comments from Christopher Hitchins, certainly an intelligent man, shortly before his death: 

``I don't have a body, I am a body." Nevertheless, he ``consciously and regularly acted as if this was not true."

I should point out that the skirmishes Hitchens engaged in between atheists and fideists are of minimal interest to us, as this post should make clear the reasons.


\hfill

\texttt{Cassiodorus on 2012-09-26 at 16:25 said: }

There was a post sometime ago that discussed a ``round table" talk that included Rene Guenon and the Thomistic philosopher Jacques Maritain. I guess I'm still beating this horse, but it seems to me that Maritain's objection to Guenon was precisely over the basic first principle of Tradition- tat tvam asi– thou art that.

If I'm not mistaken, the traditional Catholic insisted that Revelation informs metaphysics, not the other way around.

Btw, that video that was done featuring Evola and Hitchens- brilliant!!!


\hfill

\texttt{Cologero on 2012-09-26 at 20:26 said: }

Just a quick note for now regarding Maritain. In 1921 Maritain accused Guenon of participating in the rebirth of gnosis, ``mother of heresies". G. responded, ``it would just as correct to call Catholicism the father of protestantism. Actually, you are confusing gnosis and gnosticism. If you take the word gnosis in its true meaning, that of pure knowledge as I always use it, it cannot be called the mother of heresies. That is tantamount to calling truth the mother of errors."

The ``round table" was published in ``Les Nouvelles litteraires" of 25 May 1924; I don't know if that is available anywhere.

Although Maritain objected at that time to any rapprochement with Hindu metaphysics, in a later book he included Shankara along with Aristotle and Thomas as the eminent philosophers of being. Too little, too late.

BTW, you can thank Avery to that Evola-Hitchens skit.


\hfill

\texttt{Cologero on 2012-09-27 at 08:25 said: }

Regarding first principles, in ``Hindu Doctrines", Guenon explains:

\begin{quotex}
it is always questions of principle that escape the orientalists, and as it is precisely this knowledge which is essential to a proper understanding (seeing that everything else is derived from it and should logically be deduced from it), these scholars are led to neglect the one essential thing through their inability to grasp its primary importance; the consequence is that they lose their way hopelessly in a maze of the most insignificant details or in a tangle of quite arbitrary theorizing. 

\end{quotex}
This is a problem not only for orientalists, but for any ``thinker" who eschews such principles. They get mired in interminable discussions and ``debates" (at least if you include defamation as a debating tactic) and confuse a wide-ranging, though superficial, knowledge with profundity. Piling up quotes proves nothing unless the worldview or fundamental principles behind those quotes are properly addressed, whether they are implicitly inferred or explicitly admitted.


\hfill

\texttt{Anon on 2012-09-27 at 15:14 said: }

An interview with Jean-Paul Lippi : \url{http://tinyurl.com/c95wssj}

Author of « Julius Evola, métaphysicien et penseur politique. Essai d'analyse structurale (Age d'Homme, 1998) »


\end{sffamily}\end{footnotesize}
